% =============================================================
% Appendix A. Implementation Notes
% Loaded by main.tex after \appendix.
% =============================================================

\chapter{Implementation Notes: Tokenizer and Logprobs}
\label{ch:app-impl}

The probes and vocabulary statistics reported in
\S\ref{sec:mental-tokens-practice} are reproducible with a few
dozen lines of Python against a local Ollama installation and a
publicly available copy of the Gemma 4 tokenizer. This appendix
lists the moving parts and points at the scripts in the
\fname{manuscript/scripts/} directory that exercise them, in
enough detail that an interested reader can run, modify, or
port the probes themselves. No full listings; the scripts are
the listings.

\section{Inspecting Gemma 4's tokenizer}
\label{sec:app-tokenizer}

A trained tokenizer is, at the file level, a small artifact: a
flat JSON file that records everything needed to reproduce the
encoding deterministically. Sharing a tokenizer therefore
reduces to sharing this file, and most modern model releases
(Gemma 4 included) ship it alongside the weights. Google's
official Gemma weights are gated behind a license click; the
publicly mirrored copy at \code{unsloth/gemma-3-4b-it} on
HuggingFace serves the same \code{tokenizer.json}
unauthenticated, and is what produced the vocabulary statistics
in \S\ref{sec:mental-tokens-vocab}.

The JSON file holds three things: the vocabulary (a dictionary
mapping token strings to integer IDs), the ordered list of BPE
merge rules, and the set of ``added'' tokens reserved for
protocol markers (chat-template tags, padding, end-of-sequence)
that are not produced by the merge process. Token probabilities
reported elsewhere in this appendix are \emph{logprobs} ---
natural logarithms of probability, non-positive, with
$p = e^{\ell}$ recovering the probability.

The HuggingFace \code{tokenizers} library
(\code{pip install tokenizers}) loads the JSON in one call and
gives you the rest:

\begin{lstlisting}[language=Python,basicstyle=\ttfamily\footnotesize,numbers=none,frame=single,framerule=0.4pt,xleftmargin=1em,xrightmargin=1em]
from tokenizers import Tokenizer

tok = Tokenizer.from_file("tokenizer.json")
enc = tok.encode("The capital of Switzerland is Bern")
print(enc.tokens)   # ['The', 'capital', 'of', ..., 'Bern']
print(enc.ids)      # [818, 4032, 529, ..., 49870]
\end{lstlisting}

\noindent
The script \fname{manuscript/scripts/tokenization.py} wraps this
into the four operations used in the book:

\begin{itemize}[leftmargin=1.5em,itemsep=0.2em]
\item Encode a string and read off both token IDs and the
  pretty-printed pieces (\code{Tokenizer.encode}).
\item List the regular vocabulary, filtering out the added/special
  tokens reserved for chat-template markers and padding.
\item Sort by token length to find the longest entries.
\item Walk the BPE merge list, in order, to inspect the
  highest-frequency pairs at the top and the rarest at the
  bottom. The merge list itself lives in the JSON file's
  \code{model.merges} array.
\end{itemize}

\noindent
The qualitative pattern --- whitespace-only merges dominating
the first 1{,}393 entries, then English bigrams, then the long
tail into other scripts and rare strings --- is robust across
the mirrored and gated copies, and (within the limits of
cross-model comparison) across the \code{gemma4:26b} and
\code{gemma4:e2b} sibling models we use in the book.

\section{The experiments in more detail}
\label{sec:app-logprob-probes}

Once an Ollama server is reachable and a Gemma 4 variant is
pulled, the three probes of \S\ref{sec:mental-tokens-probabilities}
take three HTTP calls. The script
\fname{manuscript/scripts/next\_token\_logprobs.py} demonstrates
each as a small function; the moving parts are summarized below.

\paragraph{Base-model continuation prior (Probes 1 and 2).}
POST to \code{/api/generate} with the following body:

\begin{lstlisting}[language=json,basicstyle=\ttfamily\footnotesize,numbers=none,frame=single,framerule=0.4pt,xleftmargin=1em,xrightmargin=1em]
{
  "prompt":        "<input string>",
  "raw":           true,
  "stream":        false,
  "options":       {"num_predict": 1, "temperature": 0},
  "logprobs":      true,
  "top_logprobs":  K
}
\end{lstlisting}

\noindent
The \code{raw: true} flag suppresses Ollama's chat-template
wrapper; the model sees exactly the bytes in \code{prompt}. The
response carries a top-level \code{logprobs} array; the first
entry's \code{top\_logprobs} field is the next-token distribution
truncated to $K$ candidates. This is the call behind the
bare-prompt and trailing-space probes.

\paragraph{Chat-templated, thinking suppressed (Probe 3).}
POST to \code{/api/chat} with the following body:

\begin{lstlisting}[language=json,basicstyle=\ttfamily\footnotesize,numbers=none,frame=single,framerule=0.4pt,xleftmargin=1em,xrightmargin=1em]
{
  "model":         "<name>",
  "messages":      [{"role": "user", "content": "<question>"}],
  "stream":        false,
  "think":         false,
  "options":       {"num_predict": 1, "temperature": 0},
  "logprobs":      true,
  "top_logprobs":  K
}
\end{lstlisting}

\noindent
Two subtleties matter here. First, the OpenAI-compatible
\code{/v1/chat/completions} endpoint silently ignores
\code{think}; you must hit Ollama's native \code{/api/chat}.
Second, not every model packaging honors the flag ---
empirically, \code{gemma4:e2b} on a recent Ollama does,
\code{gemma4:26b} on an older installation does not. When the
flag is ignored, the first token of the reply is the model's
thinking-channel marker rather than the answer; the script's
demo~3 shows a hand-assembled chat-template workaround that
works regardless. The boolean in the script's profile config
controls which path is taken.

\paragraph{Profiles.} The script is parameterized by a profile
dictionary at the top of the file: model names and base URL for
each of a ``local'' (laptop's own Ollama) and a ``remote'' (LAN
Ollama with larger models) setup. Flip \code{PROFILE\_NAME} to
switch. Adding a third entry --- pointing at a different host, a
different Gemma packaging, or a non-Gemma model that exposes the
same chat/generate endpoints --- is the cleanest way to extend
the probes to a new target.
